Critical AI researchers understand the urgent need to understand and
rethink how AI systems are defined, developed, regulated, and mitigated.
\emph{Coding.Care: Guidebooks for Intersectional AI} argues that
integrating critical approaches into AI systems requires building
inviting, inclusive spaces. It requires space for more people to engage
creatively and critically with each other, and to engage creatively and
critically with machine learning as a malleable material and practice.
This project demonstrates craft-based, process-oriented approaches to AI
that can help meet this challenge.

The dissertation presents guides for fostering critical--creative coding
communities as spaces of radical belonging---activated by an ethos and a
politics modeled by queer and trans* communities that embrace radical
difference. This creates a basis for deep interdisciplinary thought,
interrogation of formative principles, and an openness to co-creation
and alternative forms necessary to reimagine AI.

With these approaches, it becomes possible to reengage emergent
technologies as craftable materials, rather than unassailable forces,
and to respond with impactful, sustainable, intersectional
interventions. Artists, activists, scholars, and technologists can
recast their relationships to emerging technologies by reframing them as
crafts like crochet---deflating AI hype, lowering barriers to learning,
and emphasizing sustainability and process. Thinking craft-as-technology
honors the inherited knowledge of many outsider communities. Thinking
technology-as-craft provides a framework to implement those theories,
ethics, and tactics as intersectional critical AI.

\emph{Coding.Care} collects four publications that enact the strategies
it theorizes. These works address and unite a wide range of communities
by relying upon different voices, forms, formats, and media. They
address various stages and processes involved, including machine
learning datasets, intersectional AI, coding communities, and embodied
algorithmic art:

\href{https://coding.care/guide.html}{Coding.Care: Field Notes for
Making Friends with Code} gives courage to pick up unfamiliar tools,
find resources to kick off a new programming project, pose questions
critically, or solve problems creatively. Its pocket guide discusses how
to build a cooperative, interdisciplinary community for co-learning
coding like the one I have facilitated since 2019. The
\href{https://intersectionalai.com}{Intersectional AI Toolkit}'s
co-authored zines are accessible guides to both AI and
intersectionality, bringing together artists, activists, academics,
makers, technologists, and anyone who wants to understand the automated
systems that impact them. The work argues that established but
marginalized tactics are necessary for reimagining more critical and
equitable machine learning.
\href{https://knowingmachines.org/critical-field-guide}{A Critical Field
Guide for Working with Machine Learning Datasets} translates critical AI
theories and data science concepts into practical tips for dataset
stewardship. Along with the
\href{https://libguides.usc.edu/inclusive-datasets}{Inclusive Datasets
Research Guide}, it combines technical skillbuilding and critical
thinking for scholars and practitioners beginning to work with large
datasets, these remain the foundation of machine learning as it grows
rapidly in impact. Lastly, a collection of five short
\href{https://coding.care}{lyric essays} offer interludes to the major
works, approaching the lived experience of an algorithmic era from
oblique angles through dialogues with five 20th century artists.

Together, the guides in \emph{Coding.Care} want to meet readers where
they are---as non-academics or those bridging into new fields, looking
for a common vocabulary to engage conscientiously with the urgent
concerns of AI systems. \emph{Coding.Care} offers an expansive
invitation to deepen interdisciplinary conversation, apply
intersectional approaches, and rework AI systems from
critical--creative--caring perspectives.
